\chapter{Экспериментальная оценка и внедрение}
\label{ch:evaluation}

\section{Методика измерений}

Требования к измерениям, принятые в работе:

\begin{enumerate}
\item Все результаты переснимаются на едином актуальном выпуске СУБД. Значения,
      приведённые в главах~\ref{ch:fanout}--\ref{ch:verification} со ссылкой на
      публикации 2017--2018~гг., относятся к выпускам того времени и служат
      исторической мотивацией, а не итоговым результатом.
\item Сравниваются сборки, различающиеся \emph{только} исследуемым изменением;
      базовая точка~--- ближайший предшествующий коммит.
\item Для каждой конфигурации выполняется не менее пяти прогонов; приводится
      медиана и разброс.
\item Одновременно с целевой метрикой измеряются метрики, которые изменение
      могло ухудшить. Для построения индекса это обязательно означает измерение
      времени выполнения запросов и перекрытия ключей~--- согласно
      §\ref{sec:overlap-fix}, ускорение построения способно замедлить поиск.
\item Указываются: версия, параметры сборки, параметры конфигурации, отличные от
      умолчаний, характеристики оборудования и файловой системы, состояние
      кэшей перед прогоном.
\end{enumerate}

% TODO: зафиксировать конкретный стенд и версию, оформить сценарии прогонов
% воспроизводимыми скриптами и приложить их к приложению.

\section{Наборы данных}

\begin{tabular}{lll}
\textbf{Набор} & \textbf{Тип} & \textbf{Назначение} \\
Равномерные точки & \code{point} & базовый случай, $D=2$ \\
Кластеризованные точки & \code{point} & чувствительность к качеству разбиения \\
Географические объекты & \code{box}, \code{polygon} & прикладной случай \\
Диапазоны дат & \code{tsrange} & односторонняя вставка, сборка мусора \\
Многомерные кубы & \code{cube} & зависимость от $D$ \\
Целочисленные массивы & \code{int[]} & обратный индекс, сигнатуры \\
Скалярные типы & \code{btree\_gist} & границы применимости сортированного построения \\
\end{tabular}

\medskip

Размерность варьируется от 2 до 32 для проверки следствий~\eqref{eq:curse};
объём~--- от помещающегося в кэш буферов до превышающего его на порядок, чтобы
разделить процессорную и вводо-выводную составляющие стоимости.

\section{Перекрытие ключей при разных стратегиях разбиения}
\label{sec:eval-overlap}

Сопоставлялись две сборки \pg{}, различающиеся ровно одним изменением
(\code{f1ea98a7975}): разбиение по заполнению страницы и разбиение функцией
\code{picksplit} на буфере в четыре страницы. Метрика~--- число обращений к
страницам индекса на точечный запрос; окно запроса масштабировалось как
$0{,}5/\sqrt{N}$, чтобы ожидаемое число совпадений не зависело от объёма.
Индекс целиком помещался в буферный кэш, поэтому измерялась работа алгоритма, а
не носителя.

Построение вставкой измерено на той же сборке: из семейства операторов удалена
опорная функция поддержки сортировки, после чего ядро выбирает обычный путь
вставки. Тем самым сравниваются алгоритмы, а не версии.

\begin{table}[h]
\centering
\caption{Три способа построения, равномерное распределение}
\begin{tabular}{llrrr}
\toprule
$N$ & способ построения & сборка, с & страниц & обращений/запрос \\
\midrule
$10^6$ & вставка & 4,38 & 9\,115 & 3,95 \\
$10^6$ & сортировка, по заполнению & 0,64 & 6\,063 & 10,01 \\
$10^6$ & сортировка, \code{picksplit} & 1,48 & 8\,165 & 4,62 \\
\midrule
$10^7$ & вставка & 69,61 & 90\,908 & 5,19 \\
$10^7$ & сортировка, по заполнению & 8,65 & 60\,608 & 12,23 \\
$10^7$ & сортировка, \code{picksplit} & 19,80 & 81\,277 & 6,10 \\
\bottomrule
\end{tabular}
\end{table}

Построение вставкой даёт дерево, близкое к оптимальному: $\omega_l$ составляет
1,0--1,13 на всех уровнях, то есть точечный запрос спускается ровно по одному
пути. Относительно этой планки разбиение по заполнению страницы удорожает
запрос в 2,5 и 2,4 раза, а разбиение функцией класса операторов~--- на 17 и
18\,\%, то есть устраняет 89 и 87\,\% штрафа. Ускорение построения относительно
вставки составляет 6,8--8,0 раза при разбиении по заполнению и 3,0--3,5 раза
при разбиении \code{picksplit}.

Отношение между двумя стратегиями сортированного построения устойчиво: 2,17 при
$10^6$ и 2,00 при $10^7$; на кластеризованных данных с размером кластера от
$0{,}5f$ до $8f$~--- 2,1--2,8.

Побочный результат, противоречащий обычному ожиданию <<лучше качество~---
больше индекс>>: сортированное построение даёт индекс \emph{меньшего} объёма,
чем построение вставкой (на 33\,\% при разбиении по заполнению и на 11\,\% при
\code{picksplit}), поскольку расщепление оставляет страницы заполненными
примерно на 70\,\%, а последовательная упаковка~--- на заданный коэффициент
заполнения.

Поуровневые значения $\omega_l$ получены восстановлением ключей снизу вверх:
для листа~--- описывающий прямоугольник его элементов, для внутреннего узла~---
объединение ключей потомков (при построении сортировкой хранимый ключ равен в
точности этому объединению).

\begin{table}[h]
\centering
\caption{Перекрытие по уровням, $N=10^7$, 2000 зондов}
\begin{tabular}{lrrrrrr}
\toprule
& \multicolumn{2}{c}{вставка} & \multicolumn{2}{c}{по заполнению} & \multicolumn{2}{c}{\code{picksplit}} \\
\cmidrule(lr){2-3}\cmidrule(lr){4-5}\cmidrule(lr){6-7}
Уровень & узлов & $\omega_l$ & узлов & $\omega_l$ & узлов & $\omega_l$ \\
\midrule
0 (корень) & 1 & 1,00 & 1 & 1,00 & 1 & 1,00 \\
1 & 8 & 1,05 & 3 & 1,78 & 5 & 1,99 \\
2 & 823 & 1,13 & 363 & 4,24 & 601 & 2,47 \\
3 (листья) & 90\,076 & 1,06 & 60\,241 & 3,95 & 80\,737 & 1,35 \\
\bottomrule
\end{tabular}
\end{table}

Величина $\omega_l$ выходит на установившееся значение и далее с глубиной не
растёт: $\omega_\infty \approx 4{,}0$ при разбиении по заполнению и
$\omega_\infty \approx 1{,}3$ при разбиении функцией класса операторов.
Верхние уровни исключение только потому, что там $\omega_l$ ограничено числом
узлов. Это подтверждает~\eqref{eq:point-steady} и объясняет устойчивость
двукратного отношения при изменении объёма на порядок.

% TODO: повторы для оценки разброса; базовая точка «построение вставкой» через
% класс операторов без опорной функции сортировки; наборы, не помещающиеся в
% буферный кэш.

\section{Построение индекса}

% TODO: время построения, объём индекса, коэффициент заполнения страниц,
% перекрытие \omega_l по уровням — для построения вставкой, наивного
% сортированного построения и сортированного построения с разбиением
% по picksplit.

\section{Поиск}

% TODO: точечный и диапазонный запрос, kNN; число прочитанных страниц против
% предсказанного \eqref{eq:point}; индексное сканирование без обращения к
% таблице при покрывающих атрибутах.

\section{Сопровождение}

% TODO: сборка мусора при логическом и физическом обходе; динамика объёма
% индекса при нагрузке «вставка в новую область, удаление из старой» с
% возвратом страниц и без; время ожидания блокировок; влияние упреждающего
% чтения.

\section{Верификация}
\label{sec:eval-verify}

\subsection{Накладные расходы}

Отношение числа прочитанных страниц к размеру индекса равно 1,00 во всех
конфигурациях (1\,134 из 1\,135, 2\,848 из 2\,849, 10\,770 из 10\,771;
расхождение~--- метастраница): обход по уровням однопроходен. Время: 0,5 / 1,3 /
4,9~с при холодном кэше и 0,0 / 0,0 / 0,1~с при прогретом, то есть проверка
ограничена вводом-выводом и на прогретом индексе почти бесплатна.

\subsection{Расход памяти и почему его нельзя измерять на одной версии}

Утверждение §\ref{sec:parent-check} состоит в том, что расход памяти ограничен
шириной одного уровня, а не размером индекса. Измерение на текущем выпуске дало
линейный рост пиковой памяти: 25,4 / 33,8 / 74,6~МБ при росте индекса в 9,5
раза и буферном кэше 32~МБ, которым этот рост не объясняется. Прочтение
напрашивалось: утверждение неверно.

Оно верно. Пара сборок вокруг \commit{1f8ab91c11e}{amcheck: Fix memory leak with
gin\_index\_check()} разделяет алгоритм и сборку:

\begin{table}[htbp]
\centering
\caption{Пиковая память процесса проверки, буферный кэш 32~МБ}
\label{tab:amcheck-mem}
\begin{tabular}{rrr}
\toprule
страниц индекса & до исправления & после исправления \\
\midrule
1\,135 & 20,1~МБ & 14,4~МБ \\
2\,849 & 28,5~МБ & 14,3~МБ \\
10\,771 & 69,8~МБ & 20,0~МБ \\
\bottomrule
\end{tabular}
\end{table}

После исправления память практически не растёт при девятикратном росте индекса,
что и утверждалось; до исправления она росла из-за утечки, независимо здесь
воспроизведённой и оценённой количественно.

Методологическое следствие существеннее самих чисел: измерение на одной версии
\emph{опровергло бы верное утверждение}, поскольку эта версия содержала дефект.
Отделить <<так устроен алгоритм>> от <<так работает данная сборка>> по одной
точке невозможно; требуется пара версий, различающихся ровно одним изменением.
Это подтверждает требование~2 методики измерений (§<<Методика измерений>>) не
как формальность, а как условие получения верного вывода.

\subsection{Полнота обнаружения}

Порча по одному инварианту с контролем посещения страницы. Достоверно
установлен один пробел: правая ссылка листовой страницы словаря, заменённая на
несуществующий номер блока, проверкой не обнаруживается, тогда как порча порядка
элементов и порча \code{pd\_lower} на \emph{той же} странице обнаруживаются обе,
то есть страница читается и разбирается.

Ещё два случая контроли перевели из <<пробел>> в <<не установлено>>: отметка
удаления при живой нисходящей ссылке (не доказано, что страница посещается) и
согласованность ключа родителя (испорчено оказалось не то, что предполагалось:
у <<последнего>> элемента внутренней страницы ключ 93 при предыдущем 19\,937,
то есть вычисленные смещения указывают не на ключ).

Отсюда требование к методике проверки полноты: молчания недостаточно. Нужно
доказать, что страница посещается, и что испорчено именно названное.

\section{Зависимость от размерности}
\label{sec:eval-dim}

Предсказание~\eqref{eq:curse} проверено на классе операторов \code{cube}~---
единственном в поставке, работающем с данными переменной размерности; для него
недоступно сортированное построение, поэтому дерево строится вставкой.
500\,000 точек, равномерно распределённых в единичном $D$-мерном кубе, 2000
зондов, оба множества порождены с фиксированным зерном. Запрос точечный: обход
посещает ровно те узлы, чей ключ содержит точку, то есть $\sum_l \omega_l$.

\begin{table}[htbp]
\centering
\caption{Число обращений к страницам при точечном запросе в зависимости от
размерности}
\label{tab:dim}
\begin{tabular}{rrrrrr}
\toprule
$D$ & $f$ & $H$ & страниц индекса & обращений & $\bar\omega$ \\
\midrule
2 & 54,8 & 4 & 4\,519 & 6,59 & 1,65 \\
3 & 38,8 & 4 & 5\,259 & 13,49 & 3,37 \\
4 & 32,4 & 4 & 6\,252 & 39,89 & 9,97 \\
6 & 17,7 & 5 & 7\,999 & 59,73 & 11,95 \\
8 & 11,6 & 5 & 9\,863 & 145,78 & 29,16 \\
12 & 5,9 & 7 & 13\,987 & 174,94 & 24,99 \\
16 & 4,0 & 8 & 18\,623 & 266,26 & 33,28 \\
24 & 2,3 & 14 & 29\,456 & 570,70 & 40,76 \\
32 & 2,0 & 16 & 38\,517 & 785,31 & 49,08 \\
\bottomrule
\end{tabular}
\end{table}

Число обращений выросло в 119 раз, но приписывать этот рост
целиком~\eqref{eq:curse} было бы неверно: за ним стоят две различные причины.

Первая~--- размер ключа: ключ \code{cube} размерности $D$ занимает $16D$ байт,
поэтому ветвистость падает с 54,8 до 2,0, а высота растёт с 4 до 16. При
ветвистости~2 дерево вырождается в двоичное, и уровней приходится проходить
больше по самому его устройству. Это следствие~\eqref{eq:fanout}, а не
геометрии.

Вторая~--- собственно перекрытие. Нормировка на высоту даёт среднее число узлов,
посещаемых на одном уровне: $\bar\omega$ растёт с 1,65 до 49,1, то есть в 30
раз. Согласно~\eqref{eq:curse} при идеальном разбиении $\omega_l \approx 1$ при
любом $D$; измеренные при $D=2$ значения 1,65 недалеки от этого, а при $D=32$ на
каждом уровне просматривается полсотни узлов вместо одного. Показатель $1/D$
делает сторону ключа близкой к единице, отклонение разбиения от идеального
перестаёт быть малым, и перекрытие растёт кратно.

Подтверждены, таким образом, направление и порядок величины, но не
функциональная форма: поуровневые значения $\omega_l$ и фактические размеры
ключей на уровнях не измерялись.

\section{Сопоставление с аналитическими моделями}

% TODO: проверка \eqref{eq:point}, \eqref{eq:build-insert}, \eqref{eq:build-sort},
% \eqref{eq:vac-logical}, \eqref{eq:vac-phys} на измеренных значениях; оценка
% области, в которой модели дают приемлемую точность.

\section{Границы применимости}

Из полученных результатов и из хода работы следуют ограничения, которые
формулируются здесь явно:

\begin{enumerate}
\item Сортированное построение требует порядка, сохраняющего локальность, и
      алгоритма разбиения, не деградирующего при четырёхкратном увеличении
      объёма входа (§\ref{sec:overlap-fix}).
\item Выигрыш от физического порядка обхода пропорционален разности стоимостей
      произвольного и последовательного чтения и сокращается на накопителях с
      малой задержкой, частично компенсируясь упреждающим чтением
      (§\ref{sec:read-stream}).
\item Возврат страниц в оборот не устраняет рост индекса полностью: внутренние
      страницы и последний потомок внутренней страницы не удаляются
      (§\ref{sec:page-deletion}).
\item При большой размерности данных перекрытие ключей растёт вне зависимости от
      качества алгоритмов~\eqref{eq:curse}; вне некоторого $D$ древовидное
      индексирование уступает последовательному просмотру.
\item Изменение функции штрафа несовместимо с требованием работоспособности
      ранее построенных индексов после обновления версии СУБД
      (§\ref{sec:penalty}).
\end{enumerate}

\section{Внедрение}

\subsection{Основная ветка СУБД PostgreSQL}

Все алгоритмы, описанные в главах~\ref{ch:fanout}--\ref{ch:verification},
приняты в основную ветку \pg{} и поставляются в составе выпусков 10--19.
Перечень изменений с идентификаторами приведён в
приложении~\ref{app:commits}.

Существенны две особенности этой формы внедрения. Во-первых, каждое изменение
прошло независимое рецензирование участниками сообщества разработчиков, и
переписка по каждому доступна публично, то есть процесс принятия результата
документирован. Во-вторых, результат распространяется вместе с СУБД, а значит
доступен всем её пользователям без дополнительных действий; это делает
эксперименты воспроизводимыми на общедоступных выпусках без обращения к
авторской реализации.

\subsection{Промышленная эксплуатация}

Результаты эксплуатируются в облачных сервисах управляемых баз данных, где
проверены на нагрузках, недостижимых в лабораторных условиях по объёму и
длительности. Именно такая эксплуатация выявила практическую значимость
результатов §\ref{sec:page-deletion} и §\ref{sec:64bit-xid}: неограниченный рост
индекса и прекращение возврата страниц в оборот проявляются на горизонте
месяцев непрерывной работы.

% TODO: акт внедрения.

\subsection{Учебный процесс}

Материалы работы используются в учебном процессе УрФУ при подготовке
бакалавров и магистров по направлению <<Информатика и вычислительная техника>>,
а также в справочной литературе по внутреннему устройству
СУБД~\cite{rogov2026}, где описаны, в частности, покрывающие атрибуты
индекса и построение обобщённого дерева поиска сортировкой.

% TODO: акт внедрения.

\subsection{Предшествующая открытая реализация}

Первой программной реализацией методов доступа, развитых в настоящей работе,
была библиотека индексирования многомерных классифицированных данных,
зарегистрированная как открытый электронный ресурс в 2009~году
\cite{imkd2009}. Распространения она не получила: в то время отсутствовала
общепринятая инфраструктура публикации исходного кода, и регистрация в
отраслевом фонде алгоритмов и программ обеспечивала формальную открытость, но
не доступность. Настоящая работа доводит тот же принцип до практического
результата: реализация распространяется не как отдельная библиотека, а как
часть СУБД, и потому доступна без каких-либо действий со стороны пользователя.

\section{Выводы по главе}

% TODO: сформулировать после получения измерений.
